\section{Introduction}
\label{sec:introduction}

In the rapidly evolving landscape of modern software development and IT operations, Infrastructure as Code (IaC) has emerged as a pivotal paradigm, fundamentally transforming how computing infrastructure is provisioned, managed and maintained. This practice defines IT infrastructure using machine-readable files, rather than relying on manual configurations on interactive tools \cite{aws_iac}. By treating infrastructure configurations as Code, IaC empowers development and operations teams to automate the deployment and management of servers, networks, databases, and other components thought version controlled configuration files \cite{chef_iac}. The adoption of IaC significantly mitigates risks associated with manual processes, such as human error and configuration drift, thereby enhancing the reliability and stability of IT systems. Furthermore, IaC contributes to substantial cost reductions, accelerates deployment cycles, and improves overall system accountability and resilience \cite{ali_iac_principles_2021, redhat_iac_2025, roper_iac_bestpractices_2025}. 

Among the various tools available for implementing IaC, Hashi-Corps Terraform stands out as a widely adopted and influential platform. This open-source IaC tool enables users to define and provision data center infrastructure using HashiCorp Configuration Language (HCL), a high-level configuration language \cite{hashicorp_terraform_iac_2025}. Its declarative nature allows users to describe the desired end-state of their infrastructure, with Terraform then automatically managing the complexities of provisioning and resource management across diverse cloud providers and no-premise environments.

While IaC and tools like Terraform have revolutionized infrastructure management, the increasing complexity and scale of modern IT environments present new challenges that traditional IaC approaches alone may struggle to address efficiently. Consequently, the integration of AI into IaC workflows has emerged as a promising avenue to further enhance automation and improve the overall development lifecycle. AI-powered IaC extends traditional capabilities by incorporating machine learning algorithms and AI-based automation into configuration management, offering potential for error detection and adherence to best practices with IaC scripts \cite{gunawat_ai-enhanced_nodate, vaiber_ai_iac_2025}. Generative AI, particularly through Large Language Models (LLMs), is increasingly being explored for automating DevOps configuration, scripting and workflow management, including generation of infrastructure code \cite{kate_generative_iac_2024}. However, despite these advancements, LLMs still face inherent limitations, such as constrains in inference-time context and a lock of structured reasoning over complex codebases \cite{khan2025kodezichronosdebuggingfirstlanguage}. Moreover, the rapid adoption of AI in code generation has raised concerns regarding code quality, intellectual property, and the potential for introducing new vulnerabilities, necessitating rigorous evaluation \cite{geladin_iac_ai_2024}.

Among the AI-powered code assistants, Github Copilot has gai-ned significant traction, offering code suggestions, agents modes, autocompletion, and snippets for various programming languages including Terraform \cite{roper_github_terraform_2024}. Its integration with tools like Terraform MCP Server further enhances its utility in IaC development \cite{thornton_why_2025, lechner_terraform_mcp_2025}. Concurrently Windsurf Cascade presents itself as an advanced AI code platform designed for multi-step code edits, autonomous feature development, and large-scale code migrations \cite{windsurf_cascade_2025}. It purports to assist in validating configurations and update through real-time feedback, combining deep codebase understanding with advanced tools \cite{windsurf_editor, saadioui_integrating_2025, sun_terraform_2025}. While both Github Copilot and Windsurf Cascade demonstrate considerable potential in augmenting IaC development, a critical gap persists in the current academic literature: a direct, empirical, and comparative evaluation of their effectiveness, efficiency, and the quality of their output specifically within the domain of IaC tasks using Terraform. Existing studies often focus on general code generation capabilities across varying levels of tasks complexity (easy, medium, complex) remains largely unexplored. This lack in systematic comparative analysis hinders a clear understanding of their practical applicability and optimal integration in real-world IaC development workflows.

This paper directly addresses the identified research gap by presenting a comprehensive empirical evaluation of Windsurf Cascade and Github Copilot in the context of Infrastructure as Code (IaC) with Terraform. Our study systematically compares the performance of these two AI-powered code assistants across a diverse set of 11 pre-defined IaC tasks, meticulously categorized by their inherent complexity into easy, medium and complex levels. To ensure the robustness and statistical significance of our findings, each task was executed three times with both AI tools, culminating in a total of 66 experimental runs. This rigorous methodological approach facilitates a direct and nuanced comparison of their efficacy, efficiency, and, crucially, the quality of the generated Terraform configurations. The insights gleaned from this research are intended to provide invaluable, data-driven understanding for both practitioners and researchers, illuminating the practical applicability and comparative advantages of Windsurf Cascade and Github Copilot within contemporary IaC development workflows. Specifically, this study contributes to the understanding how these AI tools perform under varying degrees of task complexity.